% Generated draft: requires mathematical and source review.
\documentclass[11pt]{article}
\usepackage[margin=25mm]{geometry}
\usepackage{fontspec}
\setmainfont{Latin Modern Roman}
\usepackage{amsmath,amssymb,mathtools,unicode-math}
\setmathfont{Latin Modern Math}
\usepackage{enumitem,needspace,xurl,hyperref}
\setlength{\parindent}{0pt}
\setlength{\parskip}{5pt}
\setlength{\emergencystretch}{3em}
\newcommand{\N}{\mathbb{N}}
\newcommand{\Z}{\mathbb{Z}}
\newcommand{\Q}{\mathbb{Q}}
\newcommand{\R}{\mathbb{R}}
\newcommand{\C}{\mathbb{C}}
\newcommand{\F}{\mathbb{F}}
\newcommand{\A}{\mathbb{A}}
\newcommand{\PP}{\mathbb{P}}
\newcommand{\E}{\mathbb{E}}
\newcommand{\eps}{\varepsilon}
\newcommand{\dd}{\,\mathrm{d}}
\begin{document}
\section*{UnsolvedMath ID: 701483}
\textbf{Source problem code:} catalogue entry
\subsection{Definitions and mathematical statement}
For a group \(G\), a subgroup \(H<G\) is proper if \(H\ne G\). A group \(H\) is nonabelian simple if it is nonabelian and its only normal subgroups are \(\{1\}\) and \(H\). For \(g\in G\), write \(H^g=g^{-1}Hg\).

Does there exist a group \(G\) with a proper finite nonabelian simple subgroup \(H\) such that
\[
G=\bigcup_{g\in G}H^g?
\]
Equivalently, the covering condition requires
\[
\forall x\in G\;\exists g\in G\;\exists h\in H\quad x=g^{-1}hg.
\]
No finiteness assumption is imposed on the ambient group \(G\).
\subsection{Short English statement}
Can every element of a group be conjugate, within that group, to an element of a single proper finite nonabelian simple subgroup?
\subsection{Sources}
\begin{itemize}
\item Local UnsolvedMath registry, permanent ID \(701483\), source problem code \texttt{catalogue entry}, title \emph{Kourovka Notebook Problem 18.54}; supplied upstream revision \texttt{d371ee917ac9647fbb06cd6fc529deec6f3dda82}.
\item G. Cutolo, Problem 18.54, \emph{The Kourovka Notebook (Unsolved Problems in Group Theory)}, 18th issue (2014). The supplied background also identifies the 21st issue edition (2026), arXiv:1401.0300v48. The supplied collection description names E. I. Khukhro and V. D. Mazurov as editors.
\item The supplied registry labels this problem \emph{open}. This is a supplied classification, not an independently verified assessment of its mathematical status.
\end{itemize}
\end{document}
